\section{通过变换定义点}
可用的变换包括：

\begin{itemize}
   \item 平移；
   \item 位似；
   \item 正交反射（即轴对称）；
   \item 中心对称；
   \item 正交投影；
   \item 旋转（角度制或弧度制）；
   \item 关于某个圆的反演。
\end{itemize}

\subsection{\tkzcname{tkzDefPointBy}}
变换通过选项来选择。共有两个宏：\tkzcname{tkzDefPointBy} 用于变换单个点，\tkzcname{tkzDefPointsBy} 用于变换一组点。默认情况下，点 $A$ 的像记为 $A'$。例如：
\begin{tkzltxexample}[]
\tkzDefPointBy[translation= from A to A'](B) 
\end{tkzltxexample}
结果保存在 \tkzname{tkzPointResult} 中。

\medskip
\begin{NewMacroBox}{tkzDefPointBy}{\oarg{local options}\parg{pt}}%
参数为一个已存在的点，其像保存在 \tkzname{tkzPointResult} 中。若要保留该点，可使用 \tkzcname{tkzGetPoint\{M\}} 将其命名为 \tkzname{M}。

\begin{tabular}{lll}%
\toprule
参数 &  含义 & 示例               \\ 
\midrule
\TAline{pt}   {已有点的名称}   {$(A)$}
\bottomrule
\end{tabular}

\begin{tabular}{lll}%
选项     &     & 示例                         \\ 
\midrule
\TOline{translation}{= from \#1 to \#2}{[translation=from A to B](E)}
\TOline{homothety}  {= center \#1 ratio \#2}{[homothety=center A ratio .5](E)}
\TOline{reflection} {= over \#1--\#2}{[reflection=over A--B](E)}
\TOline{symmetry }  {= center \#1}{[symmetry=center A](E)}
\TOline{projection }{= onto \#1--\#2}{[projection=onto A--B](E)}
\TOline{rotation }  {= center \#1 angle \#2}{[rotation=center O angle 30](E)}
\TOline{rotation in rad}{= center \#1 angle \#2}{[rotation in rad=center O angle pi/3](E)} 
\TOline{rotation with nodes}{= center \#1 from \#2 to \#3}{[center O from A to B](E)} 
\TOline{inversion}{= center \#1 through \#2}{[inversion =center O through A](E)} 
\TOline{inversion negative}{= center \#1 through \#2}{...} 
\bottomrule
\end{tabular}

\medskip
\emph{该宏仅定义像点，并不绘制。}
\end{NewMacroBox} 

\subsubsection{\tkzname{translation}（平移）} 

\begin{tkzexample}[latex=7cm,small]
\begin{tikzpicture}[>=latex] 
 \tkzDefPoints{0/0/A,3/1/B,3/0/C}
 \tkzDefPointBy[translation= from B to A](C) 
 \tkzGetPoint{D} 
 \tkzDrawPoints[teal](A,B,C,D)
 \tkzLabelPoints[color=teal](A,B,C,D) 
 \tkzDrawSegments[orange,->](A,B D,C) 
\end{tikzpicture} 
\end{tkzexample}

\subsubsection{\tkzname{reflection}（反射，即轴对称）} 

\begin{tkzexample}[latex=7cm,small] 
\begin{tikzpicture}[scale=.75]
 \tkzDefPoints{-2/-2/A,-1/-1/C,-4/2/D,-4/0/O}    
 \tkzDrawCircle(O,A)
 \tkzDefPointBy[reflection = over C--D](A)
 \tkzGetPoint{A'}
 \tkzDefPointBy[reflection = over C--D](O)
 \tkzGetPoint{O'}
 \tkzDrawCircle(O',A')
 \tkzDrawLine[add= .5 and .5](C,D)
 \tkzDrawPoints(C,D,O,O')
\end{tikzpicture}
\end{tkzexample}


\subsubsection{\tkzname{homothety}（位似）与 \tkzname{projection}（投影）}

\begin{tkzexample}[latex=7cm,small] 
\begin{tikzpicture}
  \tkzDefPoints{0/1/A,5/3/B,3/4/C}
  \tkzDefLine[bisector](B,A,C) \tkzGetPoint{a} 
  \tkzDrawLine[add=0 and 0,color=magenta!50 ](A,a) 
  \tkzDefPointBy[homothety=center A ratio .5](a)   
  \tkzGetPoint{a'} 
  \tkzDefPointBy[projection = onto A--B](a') 
  \tkzGetPoint{k'}  
  \tkzDefPointBy[projection = onto A--B](a) 
  \tkzGetPoint{k} 
  \tkzDrawLines[add= 0 and .3](A,k A,C)   
  \tkzDrawSegments[blue](a',k' a,k) 
  \tkzDrawPoints(a,a',k,k',A)
  \tkzDrawCircles(a',k' a,k)   
  \tkzLabelPoints(a,a',k,A)
\end{tikzpicture}
\end{tkzexample}  


\subsubsection{\tkzname{projection}（投影）}
\begin{tkzexample}[latex=7cm,small] 
\begin{tikzpicture}[scale=1.25]  
 \tkzDefPoints{0/0/A,0/4/B}
 \tkzDefTriangle[pythagore](B,A) \tkzGetPoint{C}
 \tkzDefLine[bisector](B,C,A) \tkzGetPoint{c}
 \tkzInterLL(C,c)(A,B)  \tkzGetPoint{D}
 \tkzDefPointBy[projection=onto B--C](D) 
 \tkzGetPoint{G}
 \tkzInterLC(C,D)(D,A) \tkzGetPoints{E}{F}
 \tkzDrawPolygon(A,B,C)
 \tkzDrawSegment(C,D)
 \tkzDrawCircle(D,A)
 \tkzDrawSegment[new](D,G)
 \tkzMarkRightAngle[fill=orange!10](D,G,B)
 \tkzDrawPoints(A,C,F) \tkzLabelPoints(A,C,F)
 \tkzDrawPoints(B,D,E,G)   
 \tkzLabelPoints[above right](B,D,E)
  \tkzLabelPoints[above](G)
 \end{tikzpicture}
 \end{tkzexample} 

\subsubsection{\tkzname{symmetry}（中心对称）} 
\begin{tkzexample}[latex=6cm,small] 
\begin{tikzpicture}[scale=1]
  \tkzDefPoints{2/-1/A,2/2/B,0/0/O}
  \tkzDefPointsBy[symmetry=center O](B,A){}
  \tkzDrawLine(A,A')
  \tkzDrawLine(B,B')
  \tkzMarkAngle[mark=s,arc=lll,
      size=1.5,mkcolor=red](A,O,B) 
  \tkzLabelAngle[pos=2,circle,draw,
    fill=blue!10,font=\scriptsize](A,O,B){$60^{\circ}$}
  \tkzDrawPoints(A,B,O,A',B') 
  \tkzLabelPoints(B,B')
  \tkzLabelPoints[below](A,O,A')  
\end{tikzpicture}   
\end{tkzexample}

\subsubsection{\tkzname{rotation}（旋转）} 
\begin{tkzexample}[latex=7cm,small] 
\begin{tikzpicture}[scale=0.75] 
 \tkzDefPoints{0/0/A,5/0/B}
 \tkzDrawSegment(A,B)
 \tkzDefPointBy[rotation=center A angle 60](B) 
 \tkzGetPoint{C} 
 \tkzDefPointBy[symmetry=center C](A) 
 \tkzGetPoint{D} 
 \tkzDrawSegment(A,tkzPointResult) 
 \tkzDrawLine(B,D)
 \tkzDrawArc(A,B)(C) \tkzDrawArc(B,C)(A)
 \tkzDrawArc(C,D)(D)  
 \tkzMarkRightAngle(D,B,A) 
 \tkzDrawPoints(A,B) 
 \tkzLabelPoints(A,B)
 \tkzLabelPoints[above](C)
 \tkzLabelPoints[right](D)
\end{tikzpicture}  
\end{tkzexample}  

\subsubsection{\tkzname{rotation in radian}（弧度制旋转）} 
\begin{tkzexample}[latex=6cm,small]
\begin{tikzpicture}
  \tkzDefPoint["$A$" left](1,5){A}
  \tkzDefPoint["$B$" right](4,3){B}
  \tkzDefPointBy[rotation in rad= center A angle pi/3](B)
  \tkzGetPoint{C}  
  \tkzDrawSegment(A,B)
  \tkzDrawPoints(A,B,C) 
  \tkzCompass(A,C)
  \tkzCompass(B,C) 
  \tkzLabelPoints(C)
\end{tikzpicture}
\end{tkzexample} 

\subsubsection{\tkzname{rotation with nodes}（按节点旋转）} 
\begin{tkzexample}[latex=6cm,small]
\begin{tikzpicture}
 \tkzDefPoint(0,0){O}    
 \tkzDefPoint(0:2){A} 
 \tkzDefPoint(40:2){B}  
 \tkzDefPoint(20:4){C}
 \tkzDrawLine(O,A)
 \tkzDefPointBy[rotation with nodes%
             =center O from A to B](C)  
 \tkzGetPoint{D}
\tkzDrawPoints(A,B,C,D)
\tkzDrawCircle(O,A)
\tkzLabelPoints(A,C,D)
\tkzLabelPoints[above](B)
\end{tikzpicture}
\end{tkzexample} 

\subsubsection{\tkzname{inversion}（反演）} 

反演是将点变换为其对应反演点集合的过程。设反演圆的圆心为 $O$、反演半径为 $k$，若 $P'$ 是三角形 $OQP$ 中从 $O$ 引出的高的垂足，其中 $Q$ 是圆上满足 $OQ \perp PQ$ 的一点，则称 $P$ 与 $P'$ 关于该反演圆互为反演点。\\
量 $k^2$ 称为圆幂（Coxeter 1969, p. 81）。
 
(\url{https://mathworld.wolfram.com/Inversion.html})

若干性质：
\begin{itemize}
\item 不过反演中心的圆，其反形仍为圆。
\item 过反演中心的圆，其反形为直线。
\item 不过反演中心的直线，其反形为过反演中心的圆。
\item 与反演圆正交的圆，其反形为自身。
\item 过反演中心的直线，其反形为自身。
\item 反演保持夹角不变。
\end{itemize}

说明：

直接构造：
（圆心 $O$，幂 $k^2={OA}^2=OP \times OP'$）

\begin{tkzexample}[latex=6cm,small]
\begin{tikzpicture}[scale=.5]
  \tkzDefPoints{4/0/A,6/0/P,0/0/O}
  \tkzDefPointBy[inversion = center O through A](P)
  \tkzGetPoint{P'}
  \tkzDrawSegments(O,P)
  \tkzDrawCircle(O,A)
  \tkzLabelPoints[above right,font=\scriptsize](O,A,P,P')
  \tkzDrawPoints(O,A,P,P')
\end{tikzpicture}
\end{tkzexample} 

\begin{tkzexample}[latex=6cm,small]
\begin{tikzpicture}[scale=.5]
  \tkzDefPoints{4/0/A,6/0/P,0/0/O}
  \tkzDefLine[orthogonal=through P](O,P)
  \tkzGetPoint{L}
  \tkzDefLine[tangent from = P](O,A) \tkzGetPoints{R}{Q}
  \tkzDefPointBy[projection=onto O--A](Q) \tkzGetPoint{P'}
  \tkzDrawSegments(O,P O,A)
  \tkzDrawSegments[new](O,P O,Q P,Q Q,P')
  \tkzDrawCircle(O,A)
  \tkzDrawLines[add=1 and 0](P,L)
  \tkzLabelPoints[below,font=\scriptsize](O,P')
  \tkzLabelPoints[above right,font=\scriptsize](P,Q)
  \tkzDrawPoints(O,P) \tkzDrawPoints[new](Q,P')
  \tkzLabelSegment[above](O,Q){$k$}
  \tkzMarkRightAngles(A,P',Q P,Q,O)
  \tkzLabelCircle[above=.5cm,
      font=\scriptsize](O,A)(100){inversion circle}
  \tkzLabelPoint[left,font=\scriptsize](O){inversion center}
  \tkzLabelPoint[left,font=\scriptsize](L){polar}
\end{tikzpicture}
\end{tkzexample} 


\subsubsection{直线的反演：示例 1}
\begin{tkzexample}[latex=6cm,small]  
\begin{tikzpicture}[scale=.5]
\tkzDefPoints{0/0/O,3/0/I,4/3/P,6/-3/Q}
\tkzDrawCircle(O,I)
\tkzDefPointBy[projection= onto P--Q](O) \tkzGetPoint{A}
\tkzDefPointBy[inversion = center O through I](A)
\tkzGetPoint{A'}
\tkzDefPointBy[inversion = center O through I](P)
\tkzGetPoint{P'}
\tkzDefCircle[diameter](O,A')\tkzGetPoint{o}
\tkzDrawCircle[new](o,A')
\tkzDrawLines[add=.25 and .25,red](P,Q)
\tkzDrawLines[add=.25 and .25](O,A)
\tkzDrawSegments(O,P)
\tkzDrawPoints(A,P,O) \tkzDrawPoints[new](A',P')
\end{tikzpicture}
\end{tkzexample} 

\subsubsection{直线的反演：示例 2}
\begin{tkzexample}[latex=6cm,small]  
\begin{tikzpicture}[scale=.8]
\tkzDefPoints{0/0/O,3/0/I,3/2/P,3/-2/Q}
\tkzDrawCircle(O,I)
\tkzDefPointBy[projection= onto P--Q](O) \tkzGetPoint{A}
\tkzDefPointBy[inversion = center O through I](A)
\tkzGetPoint{A'}
\tkzDefPointBy[inversion = center O through I](P)
\tkzGetPoint{P'}
\tkzDefCircle[diameter](O,A')\tkzGetPoint{o}
\tkzDrawCircle[new](o,A')
\tkzDrawLines[add=.25 and .25,red](P,Q)
\tkzDrawLines[add=.25 and .25](O,A)
\tkzDrawSegments(O,P)
\tkzDrawPoints(A,P,O) \tkzDrawPoints[new](A',P')
\end{tikzpicture}
\end{tkzexample}

\subsubsection{直线的反演：示例 3}
\begin{tkzexample}[latex=6cm,small]  
\begin{tikzpicture}[scale=.8]
\tkzDefPoints{0/0/O,3/0/I,2/1/P,2/-2/Q}
\tkzDrawCircle(O,I)
\tkzDefPointBy[projection= onto P--Q](O) \tkzGetPoint{A}
\tkzDefPointBy[inversion = center O through I](A)
\tkzGetPoint{A'}
\tkzDefPointBy[inversion = center O through I](P)
\tkzGetPoint{P'}
\tkzDefCircle[diameter](O,A')
\tkzDrawCircle[new](I,A')
\tkzDrawLines[add=.25 and .75,red](P,Q)
\tkzDrawLines[add=.25 and .25](O,A')
\tkzDrawSegments(O,P')
\tkzDrawPoints(A,P,O) \tkzDrawPoints[new](A',P')
\end{tikzpicture}
\end{tkzexample}

\subsubsection{圆的反演与 \tkzname{homothety}（位似）}
\begin{tkzexample}[latex=7cm,small]
\begin{tikzpicture}[scale=.7]
\tkzDefPoints{0/0/O,3/2/A,2/1/P}
\tkzDefLine[tangent from = O](A,P) \tkzGetPoints{T}{X}
\tkzDefPointsBy[homothety = center O
                ratio 1.25](A,P,T){}
\tkzInterCC(A,P)(A',P') \tkzGetPoints{C}{D}
\tkzCalcLength(A,P)
\tkzGetLength{rAP}
\tkzDefPointOnCircle[R=center A angle 190 radius \rAP]
\tkzGetPoint{M}
\tkzDefPointBy[inversion = center O through C](M)
\tkzGetPoint{M'}
\tkzDrawCircles[new](A,P A',P')
\tkzDrawCircle(O,C)
\tkzDrawLines[add=0 and .5](O,T' O,A' O,M' O,P')
\tkzDrawPoints(A,A',P,P',O,T,T',M,M')
\tkzLabelPoints(O,T,T',M,M')
\tkzLabelPoints[below](P,P')
\end{tikzpicture}
\end{tkzexample}


\subsubsection{三角形关于其内切圆的 \tkzname{inversion}（反演）}
\begin{tkzexample}[latex=6cm,small] 
\begin{tikzpicture}[scale=1]
\tkzDefPoints{0/0/A,5/1/B,3/6/C}
\tkzDefTriangleCenter[in](A,B,C) \tkzGetPoint{O} 
\tkzDefPointBy[projection= onto A--C](O) \tkzGetPoint{b}
\tkzDefPointBy[projection= onto A--C](O) \tkzGetPoint{b}
\tkzDefPointBy[projection= onto B--C](O) \tkzGetPoint{a}
\tkzDefPointBy[projection= onto A--B](O) \tkzGetPoint{c}
\tkzDefPointsBy[inversion = center O through b](a,b,c)%
                                             {Ia,Ib,Ic}
\tkzDefMidPoint(O,Ia) \tkzGetPoint{Ja}
\tkzDefMidPoint(O,Ib) \tkzGetPoint{Jb}
\tkzDefMidPoint(O,Ic) \tkzGetPoint{Jc}
\tkzInterCC(Ja,O)(Jb,O) \tkzGetPoints{O}{x}
\tkzInterCC(Ja,O)(Jc,O) \tkzGetPoints{y}{O}
\tkzInterCC(Jb,O)(Jc,O) \tkzGetPoints{O}{z}
\tkzDrawPolygon(A,B,C)
\tkzDrawCircle(O,b)\tkzDrawPoints(A,B,C,O)
\tkzDrawCircles[dashed,gray](Ja,y Jb,x Jc,z)
\tkzDrawArc[line width=1pt,orange,delta=0](Jb,x)(z)
\tkzDrawArc[line width=1pt,orange,delta=0](Jc,z)(y)
\tkzDrawArc[line width=1pt,orange,delta=0](Ja,y)(x)
\tkzLabelPoint[below](A){$A$}\tkzLabelPoint[above](C){$C$}
\tkzLabelPoint[right](B){$B$}\tkzLabelPoint[below](O){$O$}
\end{tikzpicture}
\end{tkzexample}

\subsubsection{与反演圆正交的圆的 \tkzname{inversion}（反演）} 
反演圆本身、与它正交的圆，以及过反演中心的直线，在反演下均保持不变。若一个圆与反演圆相交，则交点保持不变，因而这些点也位于反形圆上：见下图中的 $I$ 与 $J$。

\begin{tkzexample}[latex=5cm,small]
\begin{tikzpicture}[scale=1]
\tkzDefPoint(0,0){O}\tkzDefPoint(1,0){A}
\tkzDefPoint(-1.5,-1.5){z1} 
\tkzDefPoint(1.5,-1.25){z2} 
\tkzDefCircle[orthogonal through=z1 and z2](O,A)
\tkzGetPoint{c} 
\tkzDrawCircle[new](c,z1) 
\tkzDefPointBy[inversion =  center O through A](z1)
\tkzGetPoint{Z1} 
\tkzInterCC(O,A)(c,z1) \tkzGetPoints{I}{J}
\tkzDefPointBy[inversion =  center O through A](I)
\tkzGetPoint{I'}
\tkzDrawCircle(O,A)
\tkzDrawPoints(O,A,z1,z2) 
\tkzDrawPoints[new](c,Z1,I,J) 
\tkzLabelPoints(O,A,z1,z2,c,Z1,I,J)
\end{tikzpicture}
\end{tkzexample}



更复杂的示例参见 \tkzname{Pappus} 一节（\ref{pappus}）。

\subsubsection{\tkzname{inversion negative}（负反演）}
即先做反演，再以 $O$ 为中心做一次中心对称。
\begin{tkzexample}[latex=5cm,small]  
\begin{tikzpicture}[scale=1.5]
  \tkzDefPoints{1/0/A,0/0/O}
  \tkzDefPoint(-1.5,-1.5){z1}
  \tkzDefPoint(0.35,-2){z2} 
  \tkzDefPointBy[inversion negative = center O through A](z1)
  \tkzGetPoint{Z1} 
  \tkzDefPointBy[inversion negative = center O through A](z2)
  \tkzGetPoint{Z2}
  \tkzDrawCircle(O,A)  
  \tkzDrawPoints[color=black, fill=red,size=4](Z1,Z2)    
  \tkzDrawSegments(z1,Z1 z2,Z2)
  \tkzDrawPoints[color=black, fill=red,size=4](O,z1,z2)
  \tkzLabelPoints[font=\scriptsize](O,A,z1,z2,Z1,Z2)  
\end{tikzpicture}
\end{tkzexample} 


\newpage
\subsection{多个点的变换：\tkzcname{tkzDefPointsBy}} 
前一个宏的变体，用于同时定义多个像点。
你可以将像点的名称作为参数给出，或将该参数留空：此时像点名称由原点名称加撇号构成。

\begin{tkzltxexample}[]
\tkzDefPointsBy[translation= from A to A'](B,C){}
\end{tkzltxexample}
像点为 $B'$ 与 $C'$。

\begin{tkzltxexample}[]
\tkzDefPointsBy[translation= from A to A'](B,C){D,E} 
\end{tkzltxexample}
像点为 $D$ 与 $E$。

\begin{tkzltxexample}[]
\tkzDefPointsBy[translation= from A to A'](B)
\end{tkzltxexample}
像点为 $B'$。
\begin{NewMacroBox}{tkzDefPointsBy}{\oarg{local options}\parg{list of points}\marg{list of points}}%
\begin{tabular}{lll}%
参数 &  示例  &                  \\ 
\midrule
\TAline{\parg{list of points}\marg{list of pts}}{(A,B)\{E,F\}}{$E$、$F$ 分别为 $A$、$B$ 的像}   \\
\bottomrule
\end{tabular}

\medskip
若像点列表为空，则每个像点的名称为原点名称后加“'”。

\medskip
\begin{tabular}{lll}%
\toprule
选项     &     & 示例                         \\ 
\midrule
\TOline{translation = from \#1 to \#2}{}{[translation=from A to B](E)\{\}}
\TOline{homothety = center \#1 ratio \#2}{}{[homothety=center A ratio .5](E)\{F\}}
\TOline{reflection = over \#1--\#2}{}{[reflection=over A--B](E)\{F\}}
\TOline{symmetry = center \#1}{}{[symmetry=center A](E)\{F\}}
\TOline{projection = onto \#1--\#2}{}{[projection=onto A--B](E)\{F\}}
\TOline{rotation = center \#1 angle \#2}{}{[rotation=center  angle 30](E)\{F\}}
\TOline{rotation in rad = center \#1 angle \#2}{}{例如 angle pi/3}
\TOline{rotation with nodes = center \#1 from \#2 to \#3}{}{[center O from A to B](E)\{F\}} 
\TOline{inversion = center \#1 through \#2}{}{[inversion = center O through A](E)\{F\}} 
\TOline{inversion negative = center \#1 through \#2}{}{...} 
\bottomrule
\end{tabular}

\medskip
这些点仅被定义，并不绘制。
\end{NewMacroBox}

\subsubsection{多个点的 \tkzname{translation}（平移）}
\begin{tkzexample}[latex=7cm,small]
\begin{tikzpicture}[>=latex] 
 \tkzDefPoints{0/0/A,3/0/B,3/1/A',1/2/C}
 \tkzDefPointsBy[translation= from A to A'](B,C){} 
 \tkzDrawPolygon(A,B,C)
 \tkzDrawPolygon[new](A',B',C')
 \tkzDrawPoints(A,B,C)
 \tkzDrawPoints[new](A',B',C') 
 \tkzLabelPoints(A,B,A',B')  
 \tkzLabelPoints[above](C,C')
 \tkzDrawSegments[color = gray,->,
              style=dashed](A,A' B,B' C,C') 
\end{tikzpicture}
\end{tkzexample}

\subsubsection{多个点的 \tkzname{symmetry}（中心对称）：一个卵形}

\begin{tkzexample}[latex=7cm,small]
\begin{tikzpicture}[scale=0.4]
  \tkzDefPoint(-4,0){I}
  \tkzDefPoint(4,0){J}
  \tkzDefPoint(0,0){O} 
  \tkzInterCC(J,O)(O,J) \tkzGetPoints{L}{H}
  \tkzInterCC(I,O)(O,I) \tkzGetPoints{K}{G} 
  \tkzInterLL(I,K)(J,H) \tkzGetPoint{M}
  \tkzInterLL(I,G)(J,L) \tkzGetPoint{N}
  \tkzDefPointsBy[symmetry=center J](L,H){D,E} 
  \tkzDefPointsBy[symmetry=center I](G,K){C,F}
  \begin{scope}[line style/.style = {very thin,teal}]
    \tkzDrawLines[add=1.5 and 1.5](I,K I,G J,H J,L) 
    \tkzDrawLines[add=.5 and .5](I,J) 
    \tkzDrawCircles(O,I I,O J,O) 
    \tkzDrawArc[delta=0,orange](N,D)(C) 
    \tkzDrawArc[delta=0,orange](M,F)(E) 
    \tkzDrawArc[delta=0,orange](J,E)(D) 
    \tkzDrawArc[delta=0,orange](I,C)(F) 
  \end{scope}   
\end{tikzpicture} 
\end{tkzexample}

\endinput
